\documentclass[12pt, a4paper]{article}

\usepackage[utf8]{inputenc}
\usepackage[T1]{fontenc}
\usepackage{amsmath, amssymb, amsfonts}
\usepackage{amsthm}
\usepackage{mathtools}
\usepackage{hyperref}
\usepackage{geometry}
\usepackage{enumitem}
\usepackage{fancyhdr}
\usepackage{tikz-cd}

\geometry{margin=1in}
\setlength{\headheight}{14.5pt}
\pagestyle{fancy}
\fancyhf{}
\fancyhead[L]{Algebraic Geometry Seminar Notes}
\fancyhead[R]{UniMelb 2026 S1}
\fancyfoot[C]{\thepage}

\newtheorem{theorem}{Theorem}[section]
\newtheorem{lemma}[theorem]{Lemma}
\newtheorem{proposition}[theorem]{Proposition}
\newtheorem{corollary}[theorem]{Corollary}
\theoremstyle{definition}
\newtheorem{definition}[theorem]{Definition}
\newtheorem{example}[theorem]{Example}
\newtheorem{remark}[theorem]{Remark}

\DeclareMathOperator{\Spec}{Spec}
\DeclareMathOperator{\MaxSpec}{MaxSpec}
\DeclareMathOperator{\Hom}{Hom}
\DeclareMathOperator{\im}{im}
\DeclareMathOperator{\rad}{rad}

\newcommand{\A}{\mathbb{A}}
\newcommand{\OO}{\mathcal{O}}
\newcommand{\AffAlgSetK}{\mathrm{AffAlgSet}_K}
\newcommand{\AffVarK}{\mathrm{AffVar}_K}
\newcommand{\AlgKfgred}{\mathrm{Alg}_{K,\mathrm{fg},\mathrm{red}}}
\newcommand{\AlgKfgdom}{\mathrm{Alg}_{K,\mathrm{fg},\mathrm{dom}}}
\newcommand{\AffSch}{\mathrm{AffSch}}
\newcommand{\CRing}{\mathrm{CRing}}
\newcommand{\LRS}{\mathrm{LRS}}

\title{\textbf{Chapter 3. Coordinate Rings}}
\author{Shangren Lu}
\date{\today}

\begin{document}

\maketitle

\begin{abstract}
These notes develop affine algebraic geometry through the sequence
\[
\boxed{X \longmapsto K[X] \longmapsto \Spec(K[X]).}
\]
Starting from a classical affine algebraic set $X\subseteq \A_K^n$, we pass to its coordinate ring $K[X]$ and then to the prime spectrum. This progression distinguishes three settings: reducible classical affine geometry ($\AffAlgSetK$), the irreducible subcategory ($\AffVarK$), and the scheme-theoretic extension ($\AffSch$). Nullstellensatz identifies classical closed points with maximal ideals, whereas prime ideals parametrize irreducible closed subsets through generic points; thus the passage from maximal spectrum to prime spectrum is structurally necessary. Localization on $D(f)$ and pullback of functions provide the local and functorial mechanisms. The discussion culminates in the anti-equivalences
\[
\AffAlgSetK^{\mathrm{op}}\simeq \AlgKfgred,
\qquad
\AffSch^{\mathrm{op}}\simeq \CRing.
\]
The organizing principle is functorial: polynomial rings are free objects, quotients impose relations by universal properties, localizations invert elements universally, and affine schemes are characterized by representability statements for ring-valued morphisms.
\end{abstract}

\tableofcontents
\newpage

\section{Roadmap}

\begin{enumerate}[label=\arabic*., leftmargin=2em]
    \item Begin with a classical affine algebraic set and adopt a function-theoretic viewpoint.
    \item Justify coordinate functions via the free $K$-algebra universal property of
    \[
    K[x_1,\dots,x_n].
    \]
    \item Construct the coordinate ring by universally imposing the relations vanishing on $X$:
    \[
    K[X]\cong K[x_1,\dots,x_n]/I(X).
    \]
    \item Recover classical points via evaluation and Nullstellensatz:
    \[
    X(K)\cong \MaxSpec(K[X]).
    \]
    \item Pass from maximal ideals to prime ideals in order to represent irreducible geometry by generic points.
    \item Interpret localization through its universal property: $R_f$ is the universal $R$-algebra in which $f$ becomes invertible, corresponding to $D(f)$.
    \item Conclude with contravariance, Hom-bijections, and the two affine dictionaries
    \[
    \AffAlgSetK^{\mathrm{op}}\simeq \AlgKfgred,
    \qquad
    \AffSch^{\mathrm{op}}\simeq \CRing.
    \]
\end{enumerate}

\section{Standing conventions and terminology}

\subsection{Classical affine setting}

Throughout the classical part:
\begin{itemize}[leftmargin=2em]
    \item $K$ is an algebraically closed field;
    \item $X\subseteq \A_K^n$ is a Zariski-closed subset.
\end{itemize}

\begin{definition}[Affine algebraic set]
A subset $X\subseteq \A_K^n$ cut out by polynomial equations is called an \emph{affine algebraic set}.
\end{definition}

\begin{remark}[Terminology fix: ``variety'']
In these notes, the default classical category is $\AffAlgSetK$, allowing reducible affine algebraic sets. Its algebraic counterpart is $\AlgKfgred$, the finitely generated reduced $K$-algebras. The label ``affine variety'' is reserved for the irreducible case; one then passes to $\AffVarK$, whose algebraic counterpart is $\AlgKfgdom$.
\end{remark}

\begin{remark}[Three affine levels kept in view]
We move between three compatible settings:
\begin{itemize}[leftmargin=2em]
    \item classical reducible affine geometry: $\AffAlgSetK\leftrightarrow \AlgKfgred$;
    \item classical irreducible affine geometry: $\AffVarK\leftrightarrow \AlgKfgdom$;
    \item affine schemes over arbitrary rings: $\AffSch^{\mathrm{op}}\simeq \CRing$.
\end{itemize}
Only the first two require an algebraically closed base field; the third is the structural framework behind both.
\end{remark}

\subsection{Why category language helps here}

Category-theoretic language is used to formulate uniqueness and functoriality precisely. The basic constructions are characterized, up to unique isomorphism, by universal mapping properties: $K[x_1,\dots,x_n]$ represents $A\mapsto A^n$ on commutative $K$-algebras, $R/I$ is initial among $R$-algebras on which $I$ acts trivially, $S^{-1}R$ (in particular $R_f$) is initial among $R$-algebras in which $S$ becomes invertible, and $\Spec(R)$ is characterized by representability of the corresponding mapping functor. This viewpoint provides a uniform passage between algebraic constructions and geometric interpretations.

\section{Functions first: from equations to intrinsic algebra}

\subsection{Why functions are fundamental}

Neither an embedding $X\hookrightarrow \A_K^n$ nor a generating set for its defining ideal is canonical. The ring of polynomial functions on $X$ is intrinsic: it records polynomial observables on $X$ together with their algebraic operations. Accordingly, coordinate rings are not auxiliary data but an algebraic presentation of the geometry itself, obtained from polynomial coordinate algebras by imposing the vanishing relations.

\subsection{Polynomial algebras as free coordinate objects}

\begin{proposition}[Free commutative $K$-algebra on $n$ generators]
For every commutative $K$-algebra $A$ and every $n$-tuple $(a_1,\dots,a_n)\in A^n$, there is a unique $K$-algebra map
\[
\operatorname{ev}_{(a_1,\dots,a_n)}:K[x_1,\dots,x_n]\longrightarrow A
\]
with $x_i\mapsto a_i$ for all $i$.
Equivalently,
\[
\boxed{\Hom_{K\text{-alg}}(K[x_1,\dots,x_n],A)\cong A^n.}
\]
\end{proposition}

The commutative triangle below records this universal property; the dashed arrow is uniquely determined by $x_i\mapsto a_i$.
\[
\begin{tikzcd}[column sep=large]
K \arrow[r, "\iota"] \arrow[dr, "\eta_A"'] & K[x_1,\dots,x_n] \arrow[d, dashed, "\exists!\,\operatorname{ev}_{(a_1,\dots,a_n)}"] \\
& A
\end{tikzcd}
\]
Equivalently, $K[x_1,\dots,x_n]$ represents the functor $A\mapsto A^n$ on commutative $K$-algebras.

\begin{remark}[Naturality of polynomial coordinates]
The symbols $x_1,\dots,x_n$ serve as universal coordinate functions: every tuple in a commutative $K$-algebra is obtained from them by a unique $K$-algebra specialization. This is the precise sense in which polynomial algebras provide canonical coordinates in the affine setting.
\end{remark}

\subsection{Polynomial functions and coordinate rings}

\begin{definition}[Polynomial function on $X$]
Let $X\subseteq \A_K^n$. A \emph{polynomial function} on $X$ is a function
\[
X(K)\longrightarrow K,
\qquad
x\longmapsto f(x)
\]
obtained by restricting some $f\in K[x_1,\dots,x_n]$.
\end{definition}

\begin{definition}[Coordinate ring]
The \emph{coordinate ring} of $X$ is
\[
K[X]:=\{f|_X\, :\, f\in K[x_1,\dots,x_n]\},
\]
with pointwise ring operations.
\end{definition}

\begin{proposition}[Quotient rings impose relations universally]
Let $R$ be a commutative ring and $I\subseteq R$ an ideal. The quotient map $q:R\to R/I$ is characterized by the following universal property: for every commutative ring $A$ and every homomorphism $\varphi:R\to A$ with $I\subseteq\ker\varphi$, there exists a unique map $\overline\varphi:R/I\to A$ such that
\[
\varphi=\overline\varphi\circ q.
\]
Equivalently, $R/I$ is initial among $R$-algebras on which every element of $I$ acts by zero.
\end{proposition}

Under the condition $I\subseteq\ker\varphi$, the required factorization is represented by the commuting triangle below.
\[
\begin{tikzcd}[column sep=large]
R \arrow[r, "q"] \arrow[dr, "\varphi"'] & R/I \arrow[d, dashed, "\exists!\,\overline\varphi"] \\
& A
\end{tikzcd}
\]
Equivalently, every ring map $R\to A$ annihilating $I$ factors uniquely through $R/I$.

\begin{proposition}[Coordinate ring as universal quotient]
For $X\subseteq \A_K^n$, define
\[
I(X):=\{f\in K[x_1,\dots,x_n] : f|_X=0\}.
\]
Then there is a natural isomorphism
\[
\boxed{K[X]\cong K[x_1,\dots,x_n]/I(X).}
\]
Equivalently, $K[X]$ is obtained from the free coordinate algebra $K[x_1,\dots,x_n]$ by imposing, in a universal manner, precisely the polynomial relations vanishing on $X$.
\end{proposition}

With quotient map $q_X$ and kernel condition $\ker\rho=I(X)$, the first isomorphism theorem is encoded by the following commutative triangle.
\[
\begin{tikzcd}[column sep=large]
K[x_1,\dots,x_n] \arrow[r, "\rho"] \arrow[d, "q_X"'] & K[X] \\
K[x_1,\dots,x_n]/I(X) \arrow[ur, "\sim"'] &
\end{tikzcd}
\]
Consequently, quotienting by $I(X)$ identifies two polynomials precisely when they induce the same function on $X$.

\begin{proof}
Let
\[
\rho:K[x_1,\dots,x_n]\to K[X],\qquad f\mapsto f|_X.
\]
The map is surjective by definition of $K[X]$, and
\[
\ker\rho=\{f:f|_X=0\}=I(X).
\]
Hence, by the first isomorphism theorem,
\[
K[x_1,\dots,x_n]/I(X)\cong K[X].
\]
The universal characterization is the previous proposition applied to $R=K[x_1,\dots,x_n]$ and $I=I(X)$.
\end{proof}

\begin{corollary}[Finite type and reducedness in the classical setting]
If $X\subseteq \A_K^n$ is an affine algebraic set, then $K[X]$ is a finitely generated reduced $K$-algebra.
\end{corollary}

\begin{proof}
Finite generation follows immediately from the quotient presentation. For reducedness, let $\bar f\in K[X]$ satisfy $\bar f^m=0$ for some $m\ge 1$. Evaluating at any $x\in X(K)$ gives $\bar f(x)^m=0$ in the field $K$, hence $\bar f(x)=0$. Therefore $\bar f=0$ as a function on $X$, so $K[X]$ is reduced. Equivalently, $I(X)$ is radical; over algebraically closed $K$, this is the identity $I(V(I(X)))=\rad(I(X))$ from Nullstellensatz.
\end{proof}

\subsection{Running example: the parabola}

\begin{example}[Parabola]
Let
\[
X=V(y-x^2)\subseteq \A_K^2.
\]
Then
\[
K[X]\cong K[x,y]/(y-x^2)\cong K[x],
\]
via substitution $y=x^2$. Hence the parabola and the affine line have isomorphic coordinate rings.
\[
\boxed{K[V(y-x^2)]\cong K[x].}
\]
\end{example}

\begin{remark}[Classical notation $\Gamma(X)$]
In this classical affine setting one may write
\[
\Gamma(X)=K[X],
\]
meaning global regular functions are exactly polynomial functions on $X$.
\end{remark}

\section{Classical points, maximal ideals, and Nullstellensatz}

\subsection{Points as evaluation maps}

\begin{proposition}[Evaluation map]
For $x\in X(K)$, evaluation defines a $K$-algebra map
\[
\mathrm{ev}_x:K[X]\to K,
\qquad
f\mapsto f(x).
\]
Its kernel
\[
\mathfrak m_x:=\ker(\mathrm{ev}_x)
\]
is a maximal ideal of $K[X]$.
\end{proposition}

The relation between evaluation and residue fields is summarized by the commuting triangle below, where $\mathfrak m_x=\ker(\mathrm{ev}_x)$.
\[
\begin{tikzcd}[column sep=large]
K[X] \arrow[r, "\mathrm{ev}_x"] \arrow[d, "q_x"'] & K \\
K[X]/\mathfrak m_x \arrow[ur, "\sim"'] &
\end{tikzcd}
\]
Thus evaluation at a $K$-point agrees with passage to the residue field at the corresponding closed point.

\begin{proof}
The map is a unital $K$-algebra homomorphism and is surjective because constants lie in $K[X]$. Therefore
\[
K[X]/\mathfrak m_x\cong K,
\]
so $\mathfrak m_x$ is maximal.
\end{proof}

\begin{theorem}[Hilbert's Nullstellensatz (forms used here)]
Let $K$ be algebraically closed.
\begin{enumerate}[label=(\roman*), leftmargin=2em]
    \item For $I\subseteq K[x_1,\dots,x_n]$, one has $I(V(I))=\rad(I)$.
    \item Maximal ideals of $K[x_1,\dots,x_n]$ are exactly
    \[
    (x_1-a_1,\dots,x_n-a_n),\qquad (a_1,\dots,a_n)\in K^n.
    \]
\end{enumerate}
\end{theorem}

Applying Nullstellensatz through the quotient presentation of $K[X]$ yields the classical correspondence between points and maximal ideals. The map $x\mapsto \mathfrak m_x$ is injective because distinct points give distinct evaluation maps. For surjectivity, let $\mathfrak m\in\MaxSpec(K[X])$ and let $q:K[x_1,\dots,x_n]\twoheadrightarrow K[X]$ be the quotient map. Then $q^{-1}(\mathfrak m)$ is maximal, so
\[
q^{-1}(\mathfrak m)=(x_1-a_1,\dots,x_n-a_n)
\]
for some $(a_1,\dots,a_n)\in K^n$. The inclusion $I(X)\subseteq q^{-1}(\mathfrak m)$ implies $(a_1,\dots,a_n)\in V(I(X))$. Since $X=V(I(X))$, one has $(a_1,\dots,a_n)\in X(K)$, and hence $\mathfrak m=\mathfrak m_{(a_1,\dots,a_n)}$.
\[
\boxed{X(K)\cong \Hom_{K\text{-alg}}(K[X],K)\cong \MaxSpec(K[X]).}
\]

Consequently, one obtains the chain
\[
\boxed{X(K)\ \cong\ \MaxSpec(K[X])\ \subseteq\ \Spec(K[X]).}
\]

\begin{remark}[Jacobson viewpoint]
Finite-type $K$-algebras are Jacobson rings, so closed points (maximal ideals) remain dense and highly informative even though $\Spec$ has extra non-closed points.
\end{remark}

\begin{example}[Affine line: classical points inside the spectrum]
Let $R=K[x]$. Its maximal ideals are precisely $(x-a)$ for $a\in K$, so
\[
\MaxSpec(R)=\{(x-a):a\in K\}\cong \A^1_K(K).
\]
In $\Spec(R)$ there is one additional prime ideal, namely $(0)$, and
\[
\overline{\{(0)\}}=\Spec(K[x]).
\]
Thus $\Spec(K[x])$ consists of the classical closed points together with the non-closed generic point $(0)$.
\end{example}

\section{Why spectra are necessary: prime ideals as geometric points}

\subsection{From maximal ideals to all prime ideals}

The maximal spectrum recovers classical closed points, but irreducible closed subsets are not represented as points of $\MaxSpec$. Passing to $\Spec$ adds precisely the generic points required to encode irreducible geometry.

\begin{definition}[Prime spectrum]
For a commutative ring $R$,
\[
\Spec(R):=\{\mathfrak p\subseteq R:\ \mathfrak p\text{ prime}\}.
\]
\end{definition}

\subsection{Zariski topology}

\begin{definition}[Closed sets and standard opens]
For an ideal $I\subseteq R$, define
\[
V(I):=\{\mathfrak p\in \Spec(R): I\subseteq \mathfrak p\}.
\]
These are the closed sets. For $f\in R$, define
\[
D(f):=\Spec(R)\setminus V((f))=\{\mathfrak p\in\Spec(R):f\notin \mathfrak p\}.
\]
The sets $D(f)$ form a basis of opens.
\end{definition}

\begin{proposition}[Closed points and closure of one point]
For $\mathfrak p\in\Spec(R)$,
\[
\boxed{\overline{\{\mathfrak p\}}=V(\mathfrak p).}
\]
Consequently, $\mathfrak p$ is a closed point if and only if $\mathfrak p$ is maximal.
\end{proposition}

\begin{proof}
A closed set $V(I)$ contains $\mathfrak p$ exactly when $I\subseteq\mathfrak p$. The smallest such closed set is $V(\mathfrak p)$. It is a singleton exactly when no strictly larger prime contains $\mathfrak p$, equivalently when $\mathfrak p$ is maximal.
\end{proof}

\begin{proposition}[Irreducible closed subsets and generic points]
The assignment
\[
\mathfrak p\longmapsto V(\mathfrak p)
\]
induces a bijection
\[
\boxed{\{\text{prime ideals of }R\}\longleftrightarrow\{\text{irreducible closed subsets of }\Spec(R)\}.}
\]
Hence irreducible closed subsets are represented by prime ideals. Each irreducible closed subset has a unique generic point, namely its defining prime ideal.
\end{proposition}

\begin{remark}[Specialization and generization]
For $\mathfrak p,\mathfrak q\in\Spec(R)$, one says that $\mathfrak p$ is a specialization of $\mathfrak q$ if and only if
\[
\mathfrak q\subseteq \mathfrak p,
\]
and equivalently $\mathfrak q$ is a generization of $\mathfrak p$; this relation is precisely reverse inclusion in the poset of prime ideals.
\end{remark}

\begin{example}[Generic point of the parabola]
Let $R=K[x,y]$ and let
\[
\mathfrak p:=(y-x^2)\subseteq R.
\]
Since
\[
R/\mathfrak p\cong K[x]
\]
is a domain, $\mathfrak p$ is prime. Therefore
\[
Y:=V(\mathfrak p)\subseteq \Spec(R)
\]
is irreducible, and its generic point is exactly $\mathfrak p$. Equivalently,
\[
\overline{\{\mathfrak p\}}=V(\mathfrak p)=Y.
\]
Every closed point $\mathfrak m_{a,b}=(x-a,y-b)$ on $Y$ satisfies $b=a^2$, so
\[
\mathfrak p\subseteq (x-a,y-a^2),
\]
showing each closed point of the parabola is a specialization of the generic point.
\end{example}

\begin{example}[The affine plane in spectral language]
In $\Spec(K[x,y])$:
\begin{itemize}[leftmargin=2em]
    \item Closed points of the affine plane are the maximal ideals $(x-a,y-b)$, i.e. classical points $(a,b)\in K^2$.
    \item Each irreducible plane curve contributes a generic point: for instance $(y-3x^2)$ is the generic point of $V(y-3x^2)$, and
    \[
    \overline{\{(y-3x^2)\}}=V(y-3x^2).
    \]
    \item The ideal $(0)$ is the generic point of the whole irreducible surface, with
    \[
    \overline{\{(0)\}}=\Spec(K[x,y]).
    \]
\end{itemize}
Accordingly, spectra simultaneously encode closed points, irreducible curves via their generic points, and the generic point of the ambient irreducible space:
\[
\boxed{\text{classical points }\leftrightarrow\text{ maximal ideals }\subsetneq\text{ all prime ideals}.}
\]
\end{example}

\section{Localization and basic opens}

\subsection{Preliminaries on sheaves and affine schemes}

The scheme-theoretic viewpoint replaces a space by a pair consisting of an underlying topological space and rings of functions on its open subsets.

\begin{definition}[Sheaf of rings]
Let $X$ be a topological space. A \emph{sheaf of rings} $\mathcal F$ on $X$ consists of:
\begin{itemize}[leftmargin=2em]
    \item a ring $\mathcal F(U)$ for each open set $U\subseteq X$;
    \item a restriction map $\rho^U_V:\mathcal F(U)\to\mathcal F(V)$ for each inclusion $V\subseteq U$,
\end{itemize}
satisfying identity and composition for restrictions, together with:
\begin{enumerate}[label=(\roman*), leftmargin=2.5em]
    \item \emph{locality}: if $s,t\in\mathcal F(U)$ have equal restrictions on an open cover of $U$, then $s=t$;
    \item \emph{gluing}: compatible sections on an open cover of $U$ glue uniquely to a section of $\mathcal F(U)$.
\end{enumerate}
\end{definition}

\begin{definition}[Stalk and locally ringed space]
For $x\in X$, the \emph{stalk} $\mathcal F_x$ is the direct limit of the rings $\mathcal F(U)$ over neighborhoods $U\ni x$, equivalently the ring of germs at $x$. A \emph{locally ringed space} is a pair $(X,\OO_X)$ such that each stalk $\OO_{X,x}$ is a local ring.
\end{definition}

\begin{definition}[Affine scheme]
Let $R$ be a commutative ring and $X:=\Spec(R)$. Equipping $X$ with the structure sheaf characterized on basic opens by
\[
\OO_X(D(f))=R_f
\qquad (f\in R),
\]
one obtains the affine scheme $(X,\OO_X)=(\Spec(R),\OO_{\Spec(R)})$.
\end{definition}

\begin{example}[Punctured affine line]
Let $R=K[x]$ and $X:=\Spec(K[x])$. For the basic open
\[
D(x)=\{\mathfrak p\in\Spec(K[x]) : x\notin\mathfrak p\},
\]
one has
\[
\Gamma(D(x),\OO_X)=K[x]_x=K[x,x^{-1}].
\]
Hence $x^{-1}$ is regular on $D(x)$ but not global on $X$, since
\[
\Gamma(X,\OO_X)=K[x].
\]
This exhibits the local nature of regularity in scheme theory.
\end{example}

\subsection{Algebraic setup and notation}

Let $R$ be a commutative ring and set $X:=\Spec(R)$. For a multiplicative subset $S\subseteq R$, write
\[
i_S:R\to S^{-1}R,
\qquad
r\longmapsto \frac{r}{1},
\]
for the localization map. For $f\in R$, put
\[
S_f:=\{1,f,f^2,\dots\},
\qquad
R_f:=S_f^{-1}R,
\qquad
i_f:=i_{S_f}.
\]
For any ring $A$, the notation $A^\times$ denotes its group of units. Recall also
\[
D(f):=\{\mathfrak p\in\Spec(R): f\notin\mathfrak p\}.
\]

\begin{proposition}[Basic opens form a basis]
For $f,g\in R$,
\[
D(1)=\Spec(R),
\qquad
D(0)=\varnothing,
\qquad
D(fg)=D(f)\cap D(g).
\]
Consequently, the sets $D(f)$ form a basis of the Zariski topology on $\Spec(R)$.
\end{proposition}

\begin{proof}
The first two identities are immediate from the definitions. For the third identity,
\[
\mathfrak p\in D(fg)
\iff
fg\notin\mathfrak p
\iff
f\notin\mathfrak p\ \text{and}\ g\notin\mathfrak p
\iff
\mathfrak p\in D(f)\cap D(g),
\]
using primality of $\mathfrak p$.
\end{proof}

\subsection{Universal inversion and topology}

\begin{proposition}[Localization as universal inversion]
Let $R$ be a commutative ring and $S\subseteq R$ a multiplicative subset. The canonical map $i_S:R\to S^{-1}R$ is characterized by:
\begin{itemize}[leftmargin=2em]
    \item each $i_S(s)$ ($s\in S$) is invertible in $S^{-1}R$;
    \item for any ring map $\varphi:R\to A$ such that $\varphi(S)\subseteq A^\times$, there exists a unique map $\widetilde\varphi:S^{-1}R\to A$ with $\varphi=\widetilde\varphi\circ i_S$.
\end{itemize}
In particular, $R_f$ is the universal $R$-algebra in which $f$ becomes invertible.
\end{proposition}

For $\varphi:R\to A$ with $\varphi(f)\in A^\times$, the universal factorization is encoded by
\[
\begin{tikzcd}[column sep=large]
R \arrow[r, "i_f"] \arrow[dr, "\varphi"'] & R_f \arrow[d, dashed, "\exists!\,\widetilde\varphi"] \\
& A
\end{tikzcd}
\]

\begin{proposition}[Basic opens as localized spectra]
For $f\in R$, the map $i_f:R\to R_f$ induces a homeomorphism
\[
\Spec(R_f)\xrightarrow{\sim} D(f),
\qquad
\mathfrak q\longmapsto i_f^{-1}(\mathfrak q).
\]
Its inverse sends $\mathfrak p\in D(f)$ to the extended ideal $\mathfrak pR_f$.
\end{proposition}

\begin{proof}
Prime ideals of $R_f$ correspond to prime ideals of $R$ disjoint from $S_f$. Since $S_f=\{1,f,f^2,\dots\}$, this is equivalent to $f\notin\mathfrak p$, i.e.
\[
\mathfrak p\in D(f).
\]
Under this correspondence, contraction and extension are mutually inverse and preserve inclusions, hence induce the stated homeomorphism.
\end{proof}

\subsection{Structure sheaf on basic opens}

\begin{definition}[Structure sheaf on an affine spectrum]
The structure sheaf $\OO_X$ is the sheaf of rings determined on the basis of basic opens by
\[
\OO_X(D(f)):=R_f,
\qquad f\in R,
\]
with restriction morphisms induced by localization. For any open subset $U\subseteq X$, write $\Gamma(U,\OO_X)$ for the ring of sections on $U$.
\end{definition}

\begin{proposition}[Sections on basic opens]
For $f\in R$,
\[
\boxed{\Gamma(D(f),\OO_X)\cong R_f.}
\]
Equivalently, regular functions on $D(f)$ are exactly fractions $a/f^n$ with $a\in R$ and $n\ge 0$.
\end{proposition}

\section{From algebra back to geometry}

\subsection{Classical object-level correspondence}

\begin{theorem}[Classical affine dictionary]
Over algebraically closed $K$, there is an anti-equivalence of categories
\[
\boxed{\AffAlgSetK^{\mathrm{op}}\simeq \AlgKfgred.}
\]
If one restricts to irreducible objects, this refines to
\[
\AffVarK^{\mathrm{op}}\simeq \AlgKfgdom.
\]
\end{theorem}

\subsection{Reverse-direction example: recover equations from a ring}

\begin{example}
Let
\[
A:=K[u^2,uv,v^2]\subseteq K[u,v].
\]
Define a surjective map
\[
\psi:K[x,y,z]\twoheadrightarrow A,
\qquad
x\mapsto u^2,\ y\mapsto uv,\ z\mapsto v^2.
\]
The relation
\[
xz-y^2=0
\]
holds in $A$, and in fact
\[
\ker\psi=(xz-y^2).
\]
Hence
\[
A\cong K[x,y,z]/(xz-y^2).
\]
Therefore the associated classical affine algebraic set is
\[
V(xz-y^2)\subseteq \A_K^3,
\]
and the associated affine scheme is
\[
\Spec(A)\cong \Spec\big(K[x,y,z]/(xz-y^2)\big).
\]
\end{example}

\section{Morphisms and contravariance}

\subsection{Contravariance and pullback of functions}

A morphism $f:X\to Y$ induces, by precomposition, a homomorphism of function rings
\[
f^*:K[Y]\to K[X],\qquad g\mapsto g\circ f.
\]
Therefore the coordinate-ring construction is contravariant in the geometric variable. For schemes this yields the global-sections functor
\[
\Gamma:\AffSch^{\mathrm{op}}\longrightarrow \CRing,
\qquad
X\longmapsto \Gamma(X,\OO_X).
\]
On affine schemes, this contravariant functor is paired with $\Spec$ in the affine anti-equivalence.

\begin{definition}[Pullback on coordinate rings]
For a morphism $f:X\to Y$ of affine algebraic sets, define
\[
f^*:K[Y]\to K[X],\qquad g\mapsto g\circ f.
\]
\end{definition}

The arrow reversal is visible in the following commutative evaluation square for each $x\in X(K)$.
\[
\begin{tikzcd}[column sep=large, row sep=large]
K[Y] \arrow[r, "f^*"] \arrow[d, "\mathrm{ev}_{f(x)}"'] & K[X] \arrow[d, "\mathrm{ev}_x"] \\
K \arrow[r, equal] & K
\end{tikzcd}
\]

\begin{proposition}[A ring map induces a map of spectra]
A homomorphism $\varphi:R\to S$ induces a continuous map
\[
\Spec(S)\to\Spec(R),\qquad \mathfrak q\mapsto \varphi^{-1}(\mathfrak q).
\]
\end{proposition}

\subsection{Universal property of \texorpdfstring{$\Spec$}{Spec} and affine representability}

\begin{theorem}[Universal property of affine target]
For every locally ringed space $X$ and commutative ring $R$, there is a natural bijection
\[
\boxed{\Hom_{\LRS}(X,\Spec(R))\cong \Hom_{\CRing}(R,\Gamma(X,\OO_X)).}
\]
Equivalently, $\Spec(R)$ represents the functor
\[
X\longmapsto \Hom_{\CRing}(R,\Gamma(X,\OO_X))
\]
on locally ringed spaces.
\end{theorem}

For a ring map $\alpha:R\to\Gamma(X,\OO_X)$, the corresponding morphism $u_\alpha:X\to\Spec(R)$ is characterized by the commuting triangle below.
\[
\begin{tikzcd}[column sep=large]
R \arrow[r, "\alpha"] \arrow[d, "\cong"'] & \Gamma(X,\OO_X) \\
\Gamma(\Spec(R),\OO_{\Spec(R)}) \arrow[ur, "\Gamma(u_\alpha)"'] &
\end{tikzcd}
\]
Accordingly, a morphism into an affine scheme is determined by the induced homomorphism on global sections.

\begin{theorem}[Affine schemes Hom-bijection]
For commutative rings $R,S$,
\[
\boxed{\Hom_{\AffSch}(\Spec(S),\Spec(R))\cong \Hom_{\CRing}(R,S).}
\]
\end{theorem}

\begin{proof}
Apply the previous theorem to $X=\Spec(S)$, then use
\[
\Gamma(\Spec(S),\OO_{\Spec(S)})\cong S.
\]
\end{proof}

\begin{theorem}[Classical affine Hom-bijection]
For affine algebraic sets $X,Y$ over algebraically closed $K$,
\[
\boxed{\Hom_{\AffAlgSetK}(X,Y)\cong \Hom_{K\text{-alg}}(K[Y],K[X]).}
\]
\end{theorem}

\begin{remark}[Affine anti-equivalence]
The affine-scheme Hom-bijection yields the central structural statement:
$\Spec: \CRing^{\mathrm{op}}\to \AffSch$ and
$\Gamma: \AffSch^{\mathrm{op}}\to \CRing$ are quasi-inverse equivalences. The classical dictionary is the finite-type reduced $K$-algebra slice of this statement.
\end{remark}

\subsection{Finite limits in affine schemes}

\begin{proposition}[Products and fibre products]
Affine schemes admit finite limits; in particular, products and fibre products are given by
\[
\boxed{\Spec(R)\times_{\Spec(\mathbb Z)}\Spec(S)\cong \Spec(R\otimes_{\mathbb Z}S)}
\]
and, for $R$-algebras $A,B$,
\[
\boxed{\Spec(A)\times_{\Spec(R)}\Spec(B)\cong \Spec(A\otimes_R B).}
\]
\end{proposition}

The fibre-product formula is encoded by the pullback square below.
\[
\begin{tikzcd}[column sep=large, row sep=large]
\Spec(A\otimes_R B) \arrow[r, "\mathrm{pr}_B"] \arrow[d, "\mathrm{pr}_A"'] & \Spec(B) \arrow[d] \\
\Spec(A) \arrow[r] & \Spec(R)
\end{tikzcd}
\]
This identifies tensor product with affine base change on spectra.

\begin{proof}
Tensor products are characterized by the following mapping properties, natural in the target ring:
\[
\Hom_{\CRing}(R\otimes_{\mathbb Z}S,T)
\cong
\Hom_{\CRing}(R,T)\times \Hom_{\CRing}(S,T),
\]
\[
\Hom_{R\text{-alg}}(A\otimes_R B,T)
\cong
\Hom_{R\text{-alg}}(A,T)\times \Hom_{R\text{-alg}}(B,T).
\]
Under the affine Hom-bijection, these are exactly the universal properties of products and fibre products in $\AffSch$.
\end{proof}

\begin{remark}[Functor-of-points viewpoint]
For an affine scheme $\Spec(R)$,
\[
\boxed{h_{\Spec(R)}(A):=\Hom_{\AffSch}(\Spec(A),\Spec(R))\cong \Hom_{\CRing}(R,A).}
\]
Hence $h_{\Spec(R)}$ is representable, and the ring $R$ is recovered from its functor of points.
\end{remark}

\end{document}